% !TEX program = lualatex
% ============================================================
% 基礎数学B　前期中間試験
% ============================================================
% 使い方:
%   1. 配点を \tiny{（xx点）} の部分に記入
%   2. 問題文を各問題ブロックに記入
%   3. lualatex でコンパイル
% ============================================================

\documentclass[twocolumn,a4paper,10pt]{ltjsarticle}

% --- 日本語 ---
\usepackage{luatexja}
\usepackage[yu-win10]{luatexja-preset}

% --- レイアウト ---
\usepackage[margin=20truemm]{geometry}
\geometry{landscape}                % A4横 2段組み

% --- 数式 ---
\usepackage{amsmath,amssymb}

% --- その他 ---
\usepackage{graphicx}
\setlength{\columnseprule}{0.5pt}
\pagestyle{empty}

% ============================================================
% 編集エリア
% ============================================================
\newcommand{\ExamTitle}{基礎数学B}
\newcommand{\ExamTerm}{2026年度 前期中間試験 問題用紙}
\newcommand{\ExamClass}{（1I・1C・1A）}
\newcommand{\ExamTime}{50分}
\newcommand{\ExamTarget}{到達目標1：問題1～問題12}
% 試験範囲：第1〜7回（関数とグラフ〜２次関数と２次不等式）

% ============================================================
\begin{document}
\small
%
\begin{flushleft}
\ExamTerm \quad \ExamTitle \ExamClass 試験時間\ExamTime\\
\vspace{2mm}
学科\hspace{15mm} 出席番号 \hspace{15mm} 名前\hspace{120mm}
%
\hspace{10mm}

\begin{itemize}
\itemsep=-0.4mm
\item 各問題の解答は解答用紙の対応する枠内に記入すること．
\item 解答用紙のマークシートの学年は1を記入・マークすること．
\item 解答用紙のマークシートのクラスは，M:01，E:02，I:03，C:04，A:05に対応する2桁の数字を記入・マークすること．
\item 解答用紙のマークシートの番号は，出席番号を3桁に変換して記入・マークすること．
\item 答えしか書かれていない問題は0点となる．
\item 各問題の解答の最終的な形は，各自で適切に判断すること．
\item \ExamTarget
\end{itemize}


%
\end{flushleft}

% ============================================================
% 問題
% ============================================================

%143
\begin{description}
  \item[問題1]
  $f(x)=3x+1$のとき，$f(2)$，$f(a-2)$を求めよ．
  \tiny{（8点）}
\end{description}
\vspace{20mm}

%144(1)
\begin{description}
  \item[問題2]
  次の関数の()内の定義域に対する値域を求めよ．
  \tiny{（8点）}
  \normalsize
  \begin{equation}
    y=2x+1 \hspace{10mm} (-1 \leqq x \leqq 2)  \nonumber
  \end{equation}
\end{description}
\vspace{30mm}

%147(6)
\begin{description}
  \item[問題3]
  次の2次関数を標準形に直せ．
  \tiny{（8点）}
    \normalsize
  \begin{equation}
    y=-2x^2-3x+2  \nonumber
  \end{equation}

\end{description}
\vspace{30mm}

\newpage

%149(1)
\begin{description}
  \item[問題4]
  次の条件を満たし，$y$軸に平行な軸を持つ放物線の方程式を求めよ．
  \tiny{（8点）}
  \normalsize
  \begin{itemize}
    \item 3点$(1,0)$，$(0,3)$，$(2,-1)$を通る．
  \end{itemize}
\end{description}
\vspace{35mm}



%150(1)
\begin{description}
  \item[問題5]
  次の2次関数の最大値または最小値を求めよ．
  \tiny{（8点）}
  %
  \normalsize
  \begin{equation}
    y=x^2-4x+8  \nonumber
  \end{equation}
  %
\end{description}
\vspace{50mm}

%152
\begin{description}
  \item[問題6]
  底辺の長さと高さの和が8cmである三角形の底辺の長さを$x$cmとするとき，三角形の面積が最大となる
  $x$の値とその最大値を求めよ．
  \tiny{（8点）}
\end{description}
\vspace{50mm}

%153(1)
\begin{description}
  \item[問題7]
  次の2次関数のグラフと$x$軸との共有点を調べよ．また，共有点をもつときは，その$x$座標を求めよ．
  \tiny{（8点）}
  %
  \normalsize
  \begin{equation}
    y=2x^2+5x+1  \nonumber
  \end{equation}
  %
\end{description}
\vspace{30mm}




%154(3)
\begin{description}
  \item[問題8]
  次の2次関数のグラフが()内の条件を満たすように，定数$k$の値または値の範囲を定めよ．
  \tiny{（8点）}
  %
  \normalsize
  \begin{equation}
    y=kx^2+x-1 \hspace{10mm} (x軸と共有点をもたない) \nonumber
  \end{equation}
  %

\end{description}
\vspace{30mm}

%155(2)
\begin{description}
  \item[問題9]
  次の2次不等式を解け．
  \tiny{（8点）}
  %
  \normalsize
  \begin{equation}
    x^2-2x-1<0 \nonumber
  \end{equation}
  %
\end{description}
\vspace{30mm}

\newpage

%156改
\begin{description}
  \item[問題10]
  放物線$y=\frac{1}{2}x^2$を$x$軸方向に$-2$，$y$軸方向に1平行移動した放物線の方程式を
  求めよ．
  \tiny{（8点）}
\end{description}
\vspace{30mm}

%練習1Aの6
\begin{description}
  \item[問題11]
  2次関数$y=x^2-3x-4$のグラフを$x$軸方向にどれだけ平行移動すれば，原点を通るようにできるか．
  \tiny{（10点）}
\end{description}
\vspace{30mm}

%160
\begin{description}
  \item[問題12]
  放物線$y=-x^2+6x+1$は，放物線$y=-x^2-2x+3$をどのように平行移動したものか．
  \tiny{（10点）}
\end{description}


\end{document}
